\label{sec:threats}

\paragraph{Construct validity}
Our verdicts operationalize ``conformance,'' and a poorly chosen operationalization
could over- or under-count. Three choices limit that risk. Our checks cover
requirements that every revision in the census states in the same terms, which keeps
any server from being judged against a rule that postdates the revision it
negotiated. We also distinguish normative requirements from illustrative
ones, and the most prevalent divergence we report (unknown-tool responses) is graded as
a distinct \textsf{error-as-result} category, not as a failure, because the specification expresses it only in prose and an example, without an RFC-2119 keyword. Third, we calibrated the harness against the official reference
servers before the study and confirmed it reproduces their known-good behavior on
every check except the SDK-level divergence we set out to measure.

\paragraph{Internal validity}
The census was measured in the single-phase condition, with the network available
during both install and probe (Section~\ref{sec:method-env}). This avoids penalizing
a server that legitimately contacts the network during initialization, but it does
not remove server-side outbound calls as a confounder. The exposure is concentrated
in one check. For \textsf{tools-call-invalid-args} we accept an \texttt{isError}
result as evidence that the server rejected the ill-typed argument; an
\texttt{isError} produced instead by an unrelated runtime or network failure is
indistinguishable to the harness, and would be counted as argument checking that did
not in fact occur. The reported rate of silent type non-enforcement is therefore a
\emph{lower bound}. The lifecycle, tool-listing, error-handling, malformed-input and
stdout-purity checks are local protocol behaviors and are not exposed to this
confound. A network-isolated replication would tighten the argument-checking result,
and the harness supports that condition directly. Sandbox heterogeneity across
language runtimes is a further
source of noise, which we bound by launching each package at the version its registry
entry declares and recording every run's full launch command, including the
base-image tag.

\paragraph{External validity}
Our population is self-contained, locally-installable servers distributed on npm and
PyPI using the stdio transport. This excludes remote-only servers (which we cannot
execute without provisioning hosted infrastructure) and servers requiring
credentials (which we exclude for both ethics and reproducibility). We report the
size of each excluded class explicitly and confine all rates to the stated frame;
we do not claim they generalize to remote or credentialed servers, whose conformance
is important future work. Agent clients mostly launch local stdio servers, which is why we studied this
population first.

\paragraph{Completeness}
Because we census the eligible frame, the counts are population counts. The intervals
we report do not describe sampling error. They describe how far a rate would move if
the registry had drawn a different set of publishers, which is the uncertainty that
matters when a single publisher ships hundreds of servers from one template.

\paragraph{Publisher clustering}
The \Nprobed{} servers in our frame come from only \Npublishers{} distinct registry
namespaces, and the distribution is heavy-tailed: the largest single publisher
accounts for \TopPublisherN{} servers (\TopPublisherPct{} of the census) and the ten
largest for \TopTenPct{}. Servers from one publisher frequently share a template and
behave near-identically, so a binomial interval that treats them as independent trials
is too narrow. Every interval in this paper therefore comes from a bootstrap that
resamples publishers instead of servers. The widening is substantial: the handshake
interval is \HandshakeDesignEffect{} times wider than the binomial one
(\HandshakeRateCl{} against \HandshakeRate{}). Error-as-result becomes
\ErrAsResultRateCl{} and silent type non-enforcement \NoTypecheckRateCl{}. As a second
check we recompute each rate on a subset holding one server per namespace
(\NuniqPub{} servers): handshake yield \HandshakeRatePub{}, error-as-result
\ErrAsResultRatePub{}, silent type non-enforcement \NoTypecheckRatePub{}. No
conclusion changes under either treatment. The largest publisher is atypical in
runnability, at a \TopPublisherHandshake{} handshake yield, and its presence moves the
census figure by \TopPublisherDepressionPP{} percentage points.

\paragraph{Severity of the malformed-frame finding, and a limitation it exposed}
We initially characterized failure to recover from a malformed frame as a
denial-of-service primitive. A maintainer who reproduced the finding independently
argued that this overstates it. The argument holds. Under the stdio transport the client
launches the server and is the sole writer to its standard input, so a party capable of
sending a malformed frame can already terminate the subprocess; the specification also
places the obligation not to send invalid messages on the client. The correct reading
is a robustness and conformance failure, and this paper reports it that way. The same maintainer traced the
behavior in their server to the official Go SDK's stdio transport, so this check,
like the unknown-tool divergence, can measure an SDK default that the server only inherits.

Another maintainer's fix exposed a limitation in our instrument. In their server the
process stayed alive while its parser absorbed the frames that followed, so the next
\texttt{ping} timed out. Our harness records a single verdict for
``did not answer after a malformed frame'' and does not distinguish a terminated
process from a desynchronized one. Both are conformance failures, but they are
different failures with different fixes, and \MalformedDiesRate{} therefore aggregates
the two. Separating them requires per-server process-state tracking that the released
transcripts support but the current verdict does not encode.

\paragraph{Interrupted probes}
\(49\) probes (\(0.8\%\)) ended for reasons that may lie outside the server
(Section~\ref{sec:method-env}): \(32\) killed by the memory cap or the disk guard, and
\(17\) aborted by a harness-side error. Counting \(48\) of them as non-handshaking can
depress the handshake yield by at most \(0.8\) percentage points and does not affect
conformance rates. Each is identifiable in the released census rows by its exit code
or its missing verdicts.

\paragraph{Reproducibility}
Every number in this paper is regenerated by the released analysis scripts from the
released census rows, and we have checked that doing so reproduces them exactly. The
registry snapshot, launch commands, and package versions are recorded. Census rows do
not record the harness commit; the code that produced them is tagged
\texttt{census-v1}, and the probe logic there is unchanged from the version that was
calibrated. The registry changes daily, and the object of study is the snapshot of
19 July 2026.
A longitudinal study of the same registry shows how fast it moves, growing from 3,510
to 18,966 servers over 88.6 days~\cite{registrydrift}.
